\documentclass[11pt,twoside,openright]{book}

\usepackage[T1]{fontenc}
\usepackage[utf8]{inputenc}
\usepackage{amsmath,amssymb,amsthm,mathtools,bm}
\usepackage{mathpazo}
\usepackage{microtype}
\usepackage[inner=1.35in,outer=1.1in,top=1.25in,bottom=1.35in]{geometry}
\usepackage{booktabs}
\usepackage{fancyhdr}
\usepackage[hidelinks]{hyperref}
\usepackage{cleveref}

\setlength{\parindent}{1.4em}
\setlength{\parskip}{0pt}
\linespread{1.07}
\emergencystretch=3em

\pagestyle{fancy}
\fancyhf{}
\fancyhead[LE]{\small\thepage\quad\textsc{A Cinema of Individual Editions}}
\fancyhead[RO]{\small\nouppercase{\leftmark}\quad\thepage}
\renewcommand{\headrulewidth}{0.3pt}
\fancypagestyle{plain}{\fancyhf{}\fancyfoot[C]{\small\thepage}\renewcommand{\headrulewidth}{0pt}}

\theoremstyle{plain}
\newtheorem{theorem}{Theorem}[chapter]
\newtheorem{proposition}[theorem]{Proposition}
\newtheorem{lemma}[theorem]{Lemma}
\newtheorem{corollary}[theorem]{Corollary}
\theoremstyle{definition}
\newtheorem{definition}[theorem]{Definition}
\newtheorem{example}[theorem]{Example}
\theoremstyle{remark}
\newtheorem{remark}[theorem]{Remark}

\newcommand{\Capac}{\mathrm{Cap}}
\newcommand{\Rob}{\mathrm{Rob}}
\newcommand{\Inv}{\mathrm{Inv}}
\DeclareMathOperator{\loss}{loss}

\title{\Huge\bfseries A Cinema of Individual Editions\\[1.2ex]
\large Chapter Eleven: Locales and the Worldwide Premiere\\[0.4ex]
\large Working Draft}
\author{Flyxion\\[0.3ex]\small Independent Researcher}
\date{November 2026}

\begin{document}

\frontmatter
\maketitle
\tableofcontents
\mainmatter
\setcounter{chapter}{10}

\chapter{Locales and the Worldwide Premiere}
\label{ch:locales}

A release has one premiere instant and many local productions. Each has its own cast, language, and setting, and each serves its own audience. The promise of Chapter~1 was that no two viewers receive the same film, and the promise has a scope: within a production, which the earlier chapters have handled, or across the whole release. This chapter takes the worldwide scope. The criterion of uniqueness that Chapter~2 settled decides what composes. If uniqueness means exact inequivalence, local allocation composes when the local productions share no class. If it means separation by more than a tolerance, sharing no class is not enough, because editions selected independently in two locales must also be far enough apart across them.

The concrete case is a sequence made of shared footage, in which the local cast, language, and setting leave no perceptible trace. The chapter treats the shared specification of the story, the passages at which editions of different productions may coincide, and the minimum footage that a locale has to produce to serve its audience. Its results are computed on the lawn scene, with locales that use the pool of production A, and the footage question is also answered for an abstract pool so that its dependence on the audience and the tolerance can be seen.

\section{One Specification, Many Productions}\label{sec:spec}

What is shared among the locales is the specification of the story: the story world and its events, the script history, the scene plan with its slots and times, and the rules in force. What is local is the footage. A \emph{local production} supplies a pool of takes at each slot, shot with the locale's cast, in its language, in its setting, and the editions of a locale are drawn from its own pool.

A slot carries a \emph{local trace} if a viewer who compared the locale's takes there with those of another locale could tell them apart on account of the cast, the language, or the setting. A slot with no local trace is \emph{shared}: the same footage, or footage that differs imperceptibly, serves every locale. An establishing view of the lawn with no person in it is the standard case. Let $\ell$ be the number of slots at which every edition carries a local trace, and $s=m-\ell$ the number of shared slots. In the lawn scene the mower appears in slots $2$, $3$ and $4$, so these carry a local trace and slots $1$ and $5$, which show only the lawn, are shared.

Whether a slot carries a local trace is a fact about footage and viewers. It has to be established for each production, and a production that assumes it will have counted its way to a promise that nobody checked.

\section{Exact Inequivalence Across Locales}\label{sec:exact11}

Chapter~1 gave the condition for worldwide uniqueness under exact inequivalence. Let $G_p$ be the set of global classes that locale $p$ can realize and $U_p$ its number of viewers.

\begin{proposition}[Worldwide feasibility, restated]\label{prop:global11}
Pairwise inequivalent editions for all viewers, each viewer receiving an edition of the viewer's own locale, exist if and only if for every set $P$ of locales $\lvert\bigcup_{p\in P}G_p\rvert\ge\sum_{p\in P}U_p$. If the sets $G_p$ are pairwise disjoint, this reduces to $U_p\le\lvert G_p\rvert$ for each locale separately.
\end{proposition}

This is Proposition~1.14. Its content for the release is that allocation composes by locale exactly when the class sets are disjoint, and it fails to compose otherwise. A shared sequence is the case in which it fails to compose.

Take two locales that both use the $238$ class tuples of production A after the agreement rule, and suppose the sequence is fully shared, so that $\ell=0$ and a class tuple in one locale is the same global class as the same class tuple in the other. The union of their global classes has $238$ elements. Let the first locale have $150$ viewers and the second $100$. Each locale's audience fits within $238$, so the local test passes in both, and the combined audience of $250$ exceeds the $238$ global classes. The Hall condition fails for the pair, and some viewers in the two locales must receive the same film. If a single slot carries a local trace the global classes become $476$ and the allocation is feasible. The difference between infeasible and feasible is one slot of local footage.

\section{Perceptual Separation Across Locales}\label{sec:sep11}

If the promise is separation by more than $\varepsilon$ under the Hamming dissimilarity of Chapter~2, disjointness of classes is not enough. Two editions of different locales that differ only at slots with a local trace are different, and the question is how different.

\begin{proposition}[Cross-locale distance]\label{prop:cross}
Let every edition of every locale carry a local trace at $\ell$ slots, with local footage at a slot differing perceptibly between locales. Then any two editions of different locales differ at all $\ell$ of those slots, so their Hamming dissimilarity is at least $\ell$, and it exceeds $\ell$ by the number of shared slots at which their classes differ. Consequently, for every tolerance $\varepsilon<\ell$, locale-by-locale allocation is sound: if each locale allocates editions that are pairwise more than $\varepsilon$ apart within the locale, then all editions of the release are pairwise more than $\varepsilon$ apart, and the packing number of the union is the sum of the local packing numbers.
\end{proposition}

\begin{proof}
Editions of different locales use different footage at every local-trace slot, so they differ there, and the Hamming dissimilarity counts one for each such slot. Any pair within a locale is separated by assumption, and any cross pair has dissimilarity at least $\ell>\varepsilon$.
\end{proof}

For $\varepsilon\ge\ell$ the proposition says nothing, and the numbers show what happens. Take the lawn scene with two locales that use the pool of production A, with local traces at slots $2$, $3$ and $4$, so that $\ell=3$. \Cref{tab:sep11} gives the packing numbers. At tolerances $0$, $1$, $2$ and $3$ the union packs to exactly the sum of the local numbers. At $\varepsilon=3$ the proposition's guarantee has run out, and the sum happened to hold here because the two locales' separated sets can be chosen to differ at a shared slot. At $\varepsilon=4$ it fails: each locale packs $3$, the sum is $6$, and the union packs only $4$.

\begin{table}[ht]
\centering
\begin{tabular}{@{}lccccc@{}}
\toprule
Local-trace slots $\ell$ & $\varepsilon=0$ & $1$ & $2$ & $3$ & $4$\\
\midrule
Local packing $M_\varepsilon$ (each locale) & 238 & 54 & 16 & 7 & 3\\
Sum of the two locales & 476 & 108 & 32 & 14 & 6\\
\midrule
Union, $\ell=3$ (slots 2--4 local) & 476 & 108 & 32 & 14 & 4\\
Union, $\ell=1$ (slot 2 local) & 476 & 98 & 26 & 7 & 3\\
Union, $\ell=0$ (fully shared) & 238 & 54 & 16 & 7 & 3\\
\bottomrule
\end{tabular}
\caption{Packing numbers of the editions of two locales that use the $238$ class tuples of production A after the agreement rule, under Hamming dissimilarity, for three amounts of local trace.}\label{tab:sep11}
\end{table}

The rows with fewer local slots show the same effect earlier. With one local slot the union packs to the sum at $\varepsilon=0$ and falls short from $\varepsilon=1$, $98$ against $108$. With none, the union packs to the single locale's number at every tolerance, since the two locales' editions coincide and the second locale adds nothing.

\section{Passages Where Editions May Coincide}\label{sec:coincide}

The local-trace slots decide where coincidence can happen. Two editions of different locales coincide, as global classes, only if they agree at every slot, and they cannot agree at a slot that carries a local trace. A whole edition can therefore coincide with an edition of another locale only if the sequence has no local trace at all. A film has many scenes, though, and uniqueness can rest on any of them. A scene that is cast-free and fully shared contributes variation within a locale, and contributes nothing to separation between locales, because every locale offers the same choices there.

\begin{proposition}[Where cross-locale separation comes from]\label{prop:source}
The Hamming dissimilarity between two editions of different locales is $\ell$ plus the number of shared slots at which their classes differ. Two locales that make the same choices at every shared slot are therefore separated across locales by exactly $\ell$, and if $\ell=0$ they can coincide. Shared slots can still contribute distance when the chosen editions differ there. What a local trace supplies is a guaranteed contribution of $\ell$ irrespective of the choices at shared slots. A second distinction concerns what is counted. The packing number of the union is the size of the largest set of editions of the release that are pairwise separated. A release needs more than that: it needs a packing that meets a quota for each locale. A large combined packing may assign too many of its positions to one locale and too few to another, so the union's packing number is an upper bound on the sum of what the locales can be served and not a statement that each locale's audience is served.
\end{proposition}

The design consequence is that a production that depends on a scene being shared across locales should not count on it for worldwide distinctness. The scenes that carry cast are the scenes whose variation separates locales, and the scenes that do not are the scenes that save footage. A release plans its footage by deciding which scenes are which.

\section{The Footage a Locale Must Produce}\label{sec:footage11}

The audience of a locale is given and the footage is the cost. For the exact promise, a locale needs enough class tuples for its audience. Without further rules, a pool with $c_j$ classes at slot $j$ has $\prod_jc_j$ class tuples, and the footage is $\sum_jc_j$ classes. \Cref{tab:foot11} gives the least total number of classes over five slots that reaches an audience of $U$.

\begin{table}[ht]
\centering
\begin{tabular}{@{}rrl@{}}
\toprule
Audience $U$ & Least total classes & Classes per slot\\
\midrule
2 & 6 & $1,1,1,1,2$\\
30 & 10 & $2,2,2,2,2$\\
60 & 12 & $2,2,2,2,4$\\
238 & 15 & $3,3,3,3,3$\\
1{,}000 & 20 & $4,4,4,4,4$\\
\bottomrule
\end{tabular}
\caption{Least total number of classes over five slots whose product reaches the audience, with no rule restricting the combinations.}\label{tab:foot11}
\end{table}

The footage grows like the $m$th root of the audience. An audience of a thousand needs twenty classes, four at each slot, which is fifty times fewer classes than viewers. This is the economy the book's framing section relies on, and it holds for exact distinctness.

Tolerance changes the answer sharply. If editions must be more than $\varepsilon$ apart under the Hamming dissimilarity, a pool with $q$ classes at each of $m$ slots has at most $q^{m-\varepsilon}$ pairwise separated class tuples, by the Singleton bound, since a set of tuples that differ in more than $\varepsilon$ slots is determined by any $m-\varepsilon$ of its coordinates.

\begin{proposition}[Footage for separation]\label{prop:footsep}
To serve $U$ viewers with editions pairwise more than $\varepsilon$ apart, a uniform pool needs at least $q=\lceil U^{1/(m-\varepsilon)}\rceil$ classes at every slot. The bound is attained whenever there is a code of length $m$ over an alphabet of $q$ symbols with minimum distance $\varepsilon+1$ and at least $U$ words, for instance a Reed--Solomon code when $q$ is a prime power with $q\ge m-1$.
\end{proposition}

\begin{proof}
Two separated tuples differ in at least $\varepsilon+1$ coordinates, so their restrictions to any fixed $m-\varepsilon$ coordinates are distinct, and a set of $U$ separated tuples has $U\le q^{m-\varepsilon}$. For the attainment, a Reed--Solomon code of length $m\le q+1$ and dimension $m-\varepsilon$ has minimum distance $\varepsilon+1$ and $q^{m-\varepsilon}$ words.
\end{proof}

For five slots and an audience of sixty the bound requires $3$, $3$, $4$ and $8$ classes per slot at tolerances $0$, $1$, $2$ and $3$, against a total of $12$ classes in the best non-uniform pool at tolerance $0$, and the explicit codes reach it: a parity code over three symbols with $81$ words and minimum distance $2$ at $\varepsilon=1$, a Reed--Solomon code over a field of four elements with $64$ words and distance $3$ at $\varepsilon=2$, and one over a field of eight elements with $64$ words and distance $4$ at $\varepsilon=3$. All three were constructed and their distances checked. In total classes over five slots these are $15$, $20$ and $40$, so each step in tolerance costs more footage than the one before. For an audience of $238$ the lower bounds are $3$, $4$, $7$ and $16$, and their attainment was not checked.

Shared footage reduces what a locale must shoot. In the lawn scene let slots $1$ and $5$ be shared with two classes each, so that a locale inherits $4$ combinations from them and shoots only at the three local slots. For an audience of $30$, $60$ and $238$ the locale needs $2$, $3$ and $4$ classes at each local slot, which is $6$, $9$ and $12$ classes, against $10$, $12$ and $15$ if all five slots were local. The saving is the footage of the shared slots, and it comes at the price of \Cref{prop:source}: the locale's distinctness from other locales rests on its three local slots.

\section{What the Chapter Establishes and What It Assumes}\label{sec:limits11}

The propositions hold for any locales, any pseudometric of the form used here, and any scene plan. What the figures depend on is stated beside them: that local footage at a local-trace slot differs perceptibly between locales, which is a fact about footage and viewers; that the Hamming dissimilarity on class tuples stands in for a measured dissimilarity, as in Chapter~2; and that the scene is a five-slot lawn whose cast appears in three slots. A release with a hundred scenes has a hundred such decisions, and the chapter's counts show how the footage and the audience respond to them, not what they are.

The next chapters take the editions of a release through time. Chapter~12 asks what a ledger must retain to recover an edition on a given date, and Chapter~13 asks how a viewer's finite life limits how many of these editions the viewer can ever see.

The scripts \texttt{lawnmodel.py} and \texttt{verify\_ch11.py} reproduce each figure of this chapter.

\begin{thebibliography}{99}

\bibitem{flyxion2026gf}
Flyxion.
\newblock \emph{Generative Film: Categories, Sheaves, and Types for Cinema}.
\newblock Unpublished monograph, October 2026.

\end{thebibliography}

\end{document}
